\section{What has been established, and what has not}
\label{sec:status}

\subsection{Established}

\begin{itemize}
\item The algebraic facts of Secs.~\ref{sec:bioct} to~\ref{sec:operators}, in exact arithmetic: the
bi-octonions are alternative and not associative, their unit elements form a Moufang loop, and its
tangent algebra is a Malcev algebra that is not a Lie algebra; it contains $\mathfrak{so}(3,1)$ and eight
spinor charges with the quantum numbers of Fig.~\ref{fig:charges}.
\item Theorem R: the Malcev algebra has no non-zero realization by operators, and the loop is not a group
of operators.
\item Theorem L: no Lie algebra extends the Lorentz algebra by bosonic spinor charges whose brackets
contain Lorentz generators.
\item Theorem M: within the most general Lorentz- and isospin-covariant brackets on the sectors of the
families $\mathcal F^{\rm A}_{\rm LL}$, $\mathcal F^{\rm B}_{\rm LL}$, $\mathcal F^{\rm C}_{\rm LL}$, the
bi-octonion algebra is the unique Malcev algebra in which spinor charges close into Lorentz
generators. The necessary conditions are proved by ideal membership in Gr\"obner bases (every
polynomial used is a true consequence of the identity), and every solution is verified on all basis
tuples in exact arithmetic.
\item The manifolds of Sec.~\ref{sec:manifolds}, and the fact that the super-loop manifold is
parallelizable without a compatible Lie group structure.
\item The superfluid of Sec.~\ref{sec:superfluid} at the level of mean-field theory at zero temperature
with a sharp momentum cutoff: the channels (Theorem SF1), the spectra (SF2), the gap equation and its
critical coupling, the law of Bardeen, Cooper and Schrieffer, the ground state (SF3), its symmetry
breaking, and the super-loop singlet (SF4).
\item Power counting in $3+1$ and $7+1$ (and $4+4$) dimensions, the cutoff dependence of the contact
coupling, the divergence structure of the one-loop bubble of the Yukawa completion, the finite list of
its couplings, and the equality of loop invariance and closure invariance for quadratic forms.
\end{itemize}

\subsection{Not established}

\begin{itemize}
\item No quantum theory on a Hilbert space has the super-loop as a symmetry, and none can within
ordinary quantum mechanics (Theorem R). Whether a non-associative quantum mechanics, such as the
octonionic one of G\"unaydin, Piron and Ruegg~\cite{GunaydinPironRuegg1978}, can carry it was not
studied. The dynamical hypotheses H-dyn of the theorems on the scattering matrix are therefore not
tested.
\item The counterexample does not extend to the Poincar\'e algebra: with translations present no Malcev
algebra of our families lets spinor charges close into Lorentz generators (Theorem M(c)). Larger
families---more internal generators, central or tensorial charges, other representations of the
spinor charges---were not classified.
\item The field $\Psi_{\rm LL}$ of the superfluid is a fermion field by the spin-statistics theorem; it
carries the quantum numbers of the spinor charges, but the dynamics of Eq.~(\ref{eq:ham}) is an
ordinary relativistic field theory with the Lie symmetries $SL(2,\mathbb C)\times SU(2)_{\rm v}\times
U(1)^{\rm F}_{\rm LL}$, chosen by us. The name ``electric charge'' for $T^3_{\rm LL}$ is our choice (a
hypercharge could be added).
\item The superfluid was not treated beyond mean field: no fluctuation corrections, no finite
temperature, no dispersion of the Goldstone modes, no other interactions.
\item Renormalizability is decided by power counting and by the one-loop divergence structure of the
Yukawa completion, not by a renormalization to all orders.
\end{itemize}
